% !PRINTED 689
\weberpage{689}

The notation (5) for a differential differs from the similar notation for
the \gesperrt{functions} in $\Omega$ in that the degree of the numerator is higher by
$(2p - 2)$ than that of the denominator.

The differentials $d\alpha$ of the functions in $\Omega$ defined in § 194
% [sic] printed "definierte", the plural "Die ... Differentiale ... sind" requires "definierten"
are special cases of the general differentials (5). For not all of these
$dJ$ are differentials of functions in $\Omega$. We distinguish \gesperrt{proper
differentials} $d\alpha$, i.e. those which are differentials of functions
in $\Omega$, and \gesperrt{improper} or \gesperrt{Abelian differentials} $dJ$, which are not.
For the latter, $J$ itself has no meaning; it acquires one only in the
integral calculus, which is not accessible to the purely arithmetical
method.

If $z$ is any variable in $\Omega$ and $dJ$ a differential, then $dJ:dz$
is a function in $\alpha$,
% [sic] printed $\alpha$, a misprint for $\Omega$ (dJ:dz is a function in the field; cf. "gleichfalls immer eine Funktion in $\Omega$" two sentences later)
and $dJ:dz$ is a \gesperrt{differential quotient}. Here
too we distinguish \gesperrt{proper and improper} differential quotients. The
quotient of two differentials $dJ:dJ'$ is likewise always a function in
$\Omega$, and can likewise be called a differential quotient.

In order that a polygon quotient $\mathfrak{A}:\mathfrak{B}$ can be the
symbol for a differential, it is necessary that the difference $(a - b)$
of the orders of $\mathfrak{A}$ and $\mathfrak{B}$ be equal to $(2p - 2)$.
But this condition is not sufficient; there is the further requirement
that a variable $z$ must exist in $\Omega$ for which
\weq{6}{\mathfrak{U}^2\mathfrak{A} \sim \mathfrak{Z}\mathfrak{B}}
holds. If this requirement is satisfied, it remains satisfied when
$z$ is replaced by any other variable in $\Omega$, and $\mathfrak{A}$ and $\mathfrak{B}$ by
equivalent polygons. Thus if $\mathfrak{A}'$, $\mathfrak{A}''$, \dots{} are polygons
of the class $A$ determined by $\mathfrak{A}$, then $\mathfrak{A}':\mathfrak{B}$, $\mathfrak{A}'':\mathfrak{B}$, \dots{}
are likewise differentials.

If
\[ A = (\mathfrak{A}_1,\ \mathfrak{A}_2,\ \mathfrak{A}_3,\ \dots) \]
and $\mathfrak{A}_1$, $\mathfrak{A}_2$, $\mathfrak{A}_3$, \dots{} is a basis of $A$, then the corresponding
differential quotients
\[ \frac{dJ_1}{dz},\ \frac{dJ_2}{dz},\ \frac{dJ_3}{dz},\ \dots \]
are the basis of a family of functions of finite dimension, and we can
call
\[ (dJ_1,\ dJ_2,\ dJ_3,\ \dots) \]
the basis of a family of differentials of the same dimension.
